\documentclass[11pt]{amsart}

\usepackage[T1]{fontenc}
\usepackage[utf8]{inputenc}
\usepackage{lmodern}
\input{glyphtounicode}
\pdfgentounicode=1

\usepackage{amsmath,amssymb}
\usepackage{array,booktabs,tabularx}
\usepackage{graphicx}
\usepackage{geometry}
\usepackage{xcolor}
\usepackage{enumitem}
\usepackage{url}
\usepackage{hyperref}
\geometry{margin=1in}
\setlength{\emergencystretch}{3em}
\hypersetup{
  colorlinks=true,
  linkcolor=blue!45!black,
  citecolor=blue!45!black,
  urlcolor=blue!45!black,
  pdftitle={Adaptive Rank Allocation for Hierarchical Tensor Models via Pathwise Projected-Error Contributions},
  pdfauthor={Eliuth Chavero Jasso},
  pdfsubject={Adaptive rank allocation, hierarchical tensor networks, preregistered evaluation},
  pdfkeywords={tensor networks, rank allocation, adaptive compression, preregistration, negative results}
}

\newcommand{\norm}[1]{\lVert#1\rVert}

\theoremstyle{remark}
\newtheorem{remark}{Remark}[section]
\theoremstyle{definition}
\newtheorem{definition}{Definition}[section]

\title[Adaptive rank allocation via pathwise projected-error contributions]{Adaptive Rank Allocation for Hierarchical Tensor Models\\ via Pathwise Projected-Error Contributions}
\author{Eliuth Chavero Jasso}
\address{Independent Researcher, Apizaco, Tlaxcala, Mexico}
\date{1 August 2026}

\begin{document}
\maketitle

\begin{abstract}
We implement the primary applied benchmark motivated by the projected-error theory of
\cite{ChaveroJassoTrees}: adaptive rank allocation in a hierarchical tensor network, guided
by a pathwise score combining each node's own truncation error with an empirical estimate
of how that error amplifies along the path to the root. We evaluate this method against
five other allocation strategies and a small-case exhaustive oracle, across three
experimental levels (exact synthetic validation, teacher-student regression, and a
finite-difference Burgers-equation surrogate), with hypotheses preregistered before the
confirmatory (Level~1) run. The results are genuinely mixed: the pathwise method
significantly outperforms one of three preregistered baselines and significantly
\emph{underperforms} the other two in the confirmatory design, and is statistically
indistinguishable from most baselines in the two exploratory designs. What does hold
robustly is calibration: the pathwise majorant is an empirical upper bound in every tested
configuration, and correlates strongly with true error ($r=0.933$, $\rho=0.922$). We report
this as a negative-to-mixed result, not a success story, per the epistemic standard the
underlying mathematical work applies to itself.
\end{abstract}

\section{Introduction}

The projected-error theorem of \cite{ChaveroJassoTrees} bounds the discrepancy introduced
by orthogonal projection at every internal vertex of a finite typed multilinear composition
tree, in terms of a local closure-leakage bound $\rho$, an operator-norm bound $M$, the
number of internal vertices $k$, and a product of amplification factors along each path to
the root. This suggests a natural applied question: if a hierarchical tensor network must
allocate a limited total rank budget across its internal vertices, should that allocation be
guided by the \emph{same} local-error-times-path-amplification quantity that governs the
theorem's own bound?

We implement this as the primary applied benchmark, following the same discipline as the
mathematical work it is motivated by: every hypothesis is stated before the confirmatory
run, every comparison is reported regardless of outcome, and no numerical experiment is
treated as a proof of anything beyond what it directly measured.

\section{The network and the allocation problem}

\begin{definition}[Hierarchical tensor network, this benchmark's instantiation]
A rooted tree $T$ with leaves holding input data vectors. Each internal vertex $v$ with
children $c_1,\dots,c_a$ performs a multilinear contraction $\mu_v$ (a random or fitted
tensor of shape $(\text{ambient}_v, \dim(c_1),\dots,\dim(c_a))$), followed by an orthogonal
projection $P_v$ to a chosen rank $r_v \le \text{ambient}_v$. The projector $P_v$ is fit
once, via singular value decomposition on a batch of ambient outputs at $v$ (genuine,
data-driven singular-value decay -- not synthetic). The root is never projected.
\end{definition}

For each node we compute, exactly (Level~1) or by direct measurement (Levels~2--3):
local truncation error $\lambda_v$ (the node's own reconstruction loss, measured on
un-degraded ambient inputs, isolating it from propagated error), an empirical path
amplification factor $h_v$ (a finite-difference directional-derivative estimate of the
parent's sensitivity to $v$'s output), and the pathwise score
\[
  \mathrm{score}_v = \lambda_v \prod_{e \in \mathrm{path}(v,\mathrm{root})} h_e,
\]
exactly the mission's proposed quantity. The predicted majorant is $\sum_v \mathrm{score}_v$.

\subsection{Allocation methods compared}

Uniform; singular-energy (retain 90\% of each node's own spectral energy, rescaled to
budget); local-error-greedy (greedy on $\lambda_v$ alone); random; gradient-based (a
first-order sensitivity proxy, not full backpropagation through a continuous rank
relaxation); the proposed pathwise global-contribution method (greedy on marginal
$\mathrm{score}$ per unit rank); and a small-case exhaustive oracle. Five ablations
isolate the pathwise score's components: local-source-only, path-amplification-only, the
universal coarse $(k-1)$ bound applied as uniform allocation, a deliberately incorrect
root-residual negative control, and random path coefficients replacing the real
amplification estimates.

\section{Preregistration and experimental levels}

Hypotheses were written (\texttt{PREREGISTRATION.md}, in the accompanying software
package) before the confirmatory Level~1 campaign: primary, that pathwise allocation
achieves lower true root error than uniform, singular-energy, and local-error-greedy at
equal budget; secondary, that it reaches a fixed error tolerance at lower budget. Levels~2
and~3 were designed and run after Level~1's results were seen, and are reported as
\emph{exploratory}, not confirmatory, throughout.

\textbf{Level 1 (exact synthetic validation).} Two topologies (a depth-3 chain, a balanced
binary tree over 4 leaves), 10 seeds each, 6 rank budgets, all 6 real methods plus the
oracle plus 5 ablations: 1{,}440 records, all exactly computable (fixed random cores, no
training). \textbf{Level 2 (hierarchical tensor regression).} A teacher-student design: a
fixed full-rank ``teacher'' network defines a synthetic regression target; a ``student'' of
the same topology is fit by closed-form ridge least squares at its root, with intermediate
layers acting as fixed random multilinear features truncated per the rank allocation under
test. 5 seeds, 5 budgets, 150 records. \textbf{Level 3 (Burgers-equation surrogate).} The
same fitting scheme, with the regression target now the final state of a real
finite-difference solution of the one-dimensional viscous Burgers equation (periodic
boundary conditions, explicit upwind/central scheme) as a function of the viscosity and a
handful of initial-condition Fourier coefficients. 5 seeds, 6 budgets, 180 records.

\section{Results}

\subsection{Level 1: mixed, honestly reported}

\begin{center}
\small
\resizebox{\textwidth}{!}{%
\begin{tabular}{lrrl}
\toprule
Baseline & Mean error reduction (pathwise vs.\ baseline) & 95\% CI & Supported? \\
\midrule
uniform & $-0.055$ & $[-0.074,-0.036]$ & \textbf{No} -- uniform is better \\
singular\_energy & $+0.095$ & $[0.028,0.163]$ & \textbf{Yes} \\
local\_error\_greedy & $-0.077$ & $[-0.100,-0.057]$ & \textbf{No} -- local-greedy is better \\
\bottomrule
\end{tabular}%
}
\end{center}

This is a negative result for two of the three preregistered comparisons, retained and
reported rather than omitted. Regret against the small-case oracle confirms it:
local-error-greedy has the lowest mean regret (0.345) of all six real methods, ahead of
pathwise global-contribution (0.423). The secondary (rank-efficiency) hypothesis is
likewise not supported at any tested error tolerance.

What does hold: the pathwise majorant never exceeded the true error in any of 100 measured
configurations (ratio always in $[0.35, 0.93]$), and predicted majorant correlates strongly
with true error across all methods and configurations combined (Pearson $r=0.933$, Spearman
$\rho=0.922$, $n=1320$) -- real, positive support for calibration even where the allocation
comparison itself is not favorable.

\subsection{Ablations: informative, not merely confirmatory}

The \texttt{local\_source\_only} ablation (drop the path-amplification factor entirely)
\emph{outperforms} the full pathwise method in this design. \texttt{path\_amplification\_only}
(drop local error) is clearly worse. \texttt{random\_path\_coefficients} is statistically
indistinguishable from the real, measured path amplification (95\% CI $[-0.079,0.042]$
includes zero). Read together, these three findings say the same thing: in Level~1's
fully-random, untrained network cores, the path-amplification component of the score is not
adding value over local truncation error alone, and may be actively unhelpful. The
deliberately-incorrect root-residual negative control is confirmed worse than the real
method (CI excludes zero, correct direction) -- after a bug was found and fixed in its first
implementation (an earlier version multiplied every node's score by the same tree-wide
constant, which cannot perturb a greedy ranking at all; the corrected version uses a
genuinely node-dependent spurious factor).

\subsection{Levels 2 and 3: exploratory}

Level~2 produced a null result: all five paired comparisons against pathwise
global-contribution have 95\% confidence intervals including zero. Level~3, after a real
design bug was found and fixed (the first topology gave the one allocatable intermediate
node two one-dimensional leaf inputs, making its output mathematically rank-one regardless
of the declared ambient dimension -- every method produced byte-identical predictions at
every rank until this was corrected), gives a partially positive result: pathwise
global-contribution significantly beats uniform, local-error-greedy, and random, and ties
with singular-energy and gradient-based. The surrogate's absolute accuracy remains weak
(mean test RMSE barely below the naive training-mean baseline), a real limitation stated
plainly rather than obscured by the relative comparison.

\section{Interpretation}

Read together, the three levels do not support a universal claim that pathwise
global-contribution allocation is the best choice. What they support is narrower and more
useful: the pathwise majorant is a well-calibrated, strongly correlated estimate of true
error (a real, positive finding, criterion~3 of the mission's own success criteria), while
its use as an \emph{allocation} signal is context-dependent -- helpful against some
baselines and tasks (singular-energy in Level~1; three of five baselines in Level~3), not
others (uniform and local-error-greedy in Level~1; all baselines in Level~2). A plausible
account, not established here, is that path amplification is informative only when the
network's structure carries real long-range dependencies to exploit -- absent in Level~1's
fully random cores, present to some degree once the root is fit to real data (Levels~2--3).

\section{Limitations}

Per-node correlation between predicted contribution and each node's individually-ablated
marginal benefit (the mission's own AI5 metric, at node rather than whole-tree granularity)
was not measured. Levels~2--3 use closed-form ridge least squares at the root only, with
untrained intermediate layers -- not full end-to-end training. Level~3's absolute surrogate
accuracy is weak. No claim of a proven or certified worst-case advantage on a restricted
class (mission success criterion~4) was attempted.

\section{Reproducibility}

All code, raw records (1{,}770 total across three levels), and analysis scripts are in
\texttt{applications/adaptive\_tensor\_network/} in the accompanying repository. Every
result in this paper is read directly from a committed JSON file produced by a script also
in that directory; none is asserted from memory. 16 unit/regression tests cover budget
conservation for every method and ablation, projector orthonormality and idempotence, no
test-set leakage, determinism given a fixed seed, and both bugs described above (as
named regression tests, so they cannot silently reappear).

\begin{thebibliography}{9}
\bibitem{ChaveroJassoTrees}
E.~Chavero Jasso, \emph{Recursive Orthogonal Projection of Typed Multilinear Trees},
accompanying manuscript, this package.
\end{thebibliography}

\end{document}
